\chapter{An\'alisis de antecedentes y aportaci\'on realizada}\label{analanteced}

\section{Introducci\'on}

Este cap\'itulo analiza la situaci\'on de partida del dominio de gesti\'on de entregables acad\'emicos
y describe con detalle la aportaci\'on concreta del proyecto sobre dicho contexto.
El an\'alisis se sustenta en los artefactos de requisitos y dise\~no (ERS y DAS),
y se contrasta con el estado real de implementaci\'on reflejado en el backlog,
el plan de sprints y la estrategia de pruebas.

El objetivo de esta secci\'on no es solo comparar un ``antes'' y un ``despu\'es'',
sino explicar por qu\'e el problema exist\'ia,
qu\'e brechas de ingenier\'ia lo sosten\'ian,
y c\'omo la soluci\'on propuesta aporta una mejora estructural defendible.

\section{Metodolog\'ia de an\'alisis de antecedentes}

Se adopta un enfoque en tres pasos: la caracterizaci\'on del escenario \emph{actual} (actores, procesos, fortalezas y debilidades), la identificaci\'on de brechas operativas y tecnol\'ogicas, y la definici\'on del escenario \emph{objetivo} con mapeo expl\'icito de aportaciones. Este orden evita saltos l\'ogicos y permite justificar cada mejora a partir de una necesidad concreta.

Este m\'etodo evita descripciones gen\'ericas y permite mantener trazabilidad
entre problema detectado, decisi\'on de dise\~no e impacto obtenido.

\section{Situaci\'on actual (escenario \emph{actual})}

\subsection{Actores y procesos principales}

En la situaci\'on de partida intervienen tres actores de negocio: el \textbf{profesor}, que crea actividades y gestiona la recogida y correcci\'on de entregas; el \textbf{estudiante}, que realiza la entrega de ficheros en la plataforma docente; y el \textbf{administrador de plataforma}, responsable de mantener servicios base e infraestructura institucional. Esta distribuci\'on delimita responsabilidades y explica por qu\'e ciertos problemas se concentran en el flujo docente.

Los procesos de referencia en ese escenario son \textbf{PRON-001} (creaci\'on y publicaci\'on de tareas), \textbf{PRON-002} (entrega de pr\'acticas por el alumnado) y \textbf{PRON-003} (descarga masiva y organizaci\'on manual de entregas por el profesorado). La presencia de un proceso dedicado a organizaci\'on manual evidencia el principal foco de fricci\'on operativa.

\subsection{Fortalezas del entorno de partida}

El an\'alisis del contexto inicial evidencia fortalezas relevantes. En particular, \textbf{FOR-001} aporta centralizaci\'on de usuarios mediante identidad institucional, \textbf{FOR-002} garantiza disponibilidad generalizada y acceso web estandarizado, y \textbf{FOR-003} ofrece un marco de uso conocido por estudiantes y docentes. Estas ventajas justifican una estrategia de complementariedad en lugar de sustituci\'on total.

Estas fortalezas justifican que la propuesta no persiga sustituir por completo
el ecosistema acad\'emico existente,
sino complementarlo en el punto donde se concentra la fricci\'on operativa.

\subsection{Debilidades estructurales detectadas}

Junto a las fortalezas, se detectan debilidades que afectan al rendimiento del proceso. La \textbf{DEB-001} se manifiesta en rigidez de gesti\'on de archivos y clasificaci\'on manual reiterada, mientras que la \textbf{DEB-002} limita la evoluci\'on funcional al depender de arquitecturas cerradas. A esto se suma la \textbf{DEB-003}, que reduce la flexibilidad del flujo de entrega para el estudiante, y la \textbf{DEB-004}, que refleja una integraci\'on no nativa con almacenamiento cloud moderno. En conjunto, estas debilidades concentran el coste operativo en fases cr\'iticas de entrega y correcci\'on.

\subsection{An\'alisis de causas ra\'iz}

Las debilidades anteriores no se explican por errores puntuales, sino por causas sist\'emicas. En el plano de proceso predomina la carga manual en fases cr\'iticas de correcci\'on, y en el plano de producto falta especializaci\'on para el ciclo completo de entregables. A nivel t\'ecnico existe acoplamiento entre flujo acad\'emico y herramientas de prop\'osito general, y en calidad se observa una trazabilidad escasa entre entrega, versi\'on, evaluaci\'on y feedback. Esta combinaci\'on explica por qu\'e las mejoras parciales no eliminan la fricci\'on operativa.

\section{Brechas entre escenario actual y necesidades objetivo}

Para justificar la intervenci\'on,
se identifican brechas concretas entre la situaci\'on \emph{actual}
y el resultado esperado en t\'erminos de ingenier\'ia.

\begin{table}[ht]
\centering
\begin{tabular}{|P{3.2cm}|P{4.4cm}|P{4.9cm}|}
\hline
\textbf{Brecha} & \textbf{Situaci\'on inicial} & \textbf{Necesidad objetivo} \\
\hline
Operativa docente & Descarga y organizaci\'on manual recurrente & Flujo integrado de correcci\'on con menor carga administrativa \\
\hline
Experiencia de entrega & Interacci\'on poco guiada seg\'un contexto & Flujo orientado por rol y estado de entregable \\
\hline
Trazabilidad acad\'emica & Datos dispersos entre herramientas y estados & Cadena unificada entrega-version-evaluaci\'on-feedback \\
\hline
Evoluci\'on funcional & Alta dependencia de plataforma generalista & Arquitectura modular y extensible por componentes \\
\hline
Calidad y despliegue & Validaci\'on heterog\'enea por tarea & Verificaci\'on automatizada y entrega continua \\
\hline
\end{tabular}
\caption{Brechas identificadas entre el escenario de partida y el objetivo}
\label{tab:brechas-antecedentes}
\end{table}

\section{Entorno tecnol\'ogico de partida}

El contexto previo se caracteriza por una plataforma institucional LMS con autenticaci\'on centralizada, acceso multiplataforma desde navegador y dependencia de herramientas locales para el tratamiento de paquetes de entrega. Esta dependencia traslada parte del esfuerzo al puesto de trabajo del docente y amplifica diferencias de criterio en la gesti\'on de materiales.

La consecuencia directa es que parte del trabajo cr\'itico de evaluaci\'on
se desplaza al puesto local del docente,
con impacto en tiempo de correcci\'on,
estandarizaci\'on de criterios
y riesgo de inconsistencias en trazabilidad.

\section{An\'alisis de competidores y brechas sectoriales}

Para dimensionar la aportaci\'on real del proyecto en el sector,
se revis\'o el comportamiento de soluciones ampliamente usadas
en docencia, as\'i como herramientas m\'as espec\'ificas de recepci\'on
de entregas. El objetivo no es enumerar marcas, sino identificar
patrones funcionales y brechas que condicionan la operativa docente.

\subsection{Tipolog\'ias de competidores analizados}

El an\'alisis se ha estructurado en cuatro familias. Se consideran \textbf{LMS generalistas} que priorizan gesti\'on de cursos, contenidos, evaluaciones y comunicaci\'on, \textbf{herramientas de evaluaci\'on t\'ecnica} orientadas a entregas de programaci\'on o autoevaluaci\'on, \textbf{suites colaborativas y almacenamiento cloud} centradas en compartir archivos y carpetas sin cubrir el flujo completo de entrega, y \textbf{gestores de entrega simples} que ofrecen subida de ficheros con m\'inimo control posterior. Esta segmentaci\'on permite comparar la propuesta con alternativas reales sin reducir el an\'alisis a marcas concretas.

\subsection{LMS generalistas}

Las plataformas LMS ofrecen una cobertura funcional amplia y estable,
pero su dise\~no generalista tiende a tratar las entregas como un m\'odulo
m\'as de un ecosistema mayor. Esto provoca que la validaci\'on de ficheros
sea habitualmente limitada a extensi\'on y tama\~no, y que el control
de estructura interna de paquetes ZIP recaiga en instrucciones manuales
al alumnado o en revisi\'on posterior del docente.

\subsection{Herramientas de evaluaci\'on t\'ecnica}

Las plataformas especializadas en correcci\'on t\'ecnica priorizan
integraciones con repositorios o sistemas de autoevaluaci\'on.
Son muy potentes cuando la entrega es estrictamente c\'odigo,
pero pierden flexibilidad cuando el entregable es heterog\'eneo
(documentos, material multimedia, informes mixtos o combinaciones).

\subsection{Suites colaborativas y almacenamiento cloud}

Las suites colaborativas permiten compartir archivos y carpetas
de forma flexible, pero no sustituyen un flujo acad\'emico completo.
En la pr\'actica, la trazabilidad de versiones, los estados de entrega
y la publicaci\'on de notas suelen quedar fuera del sistema,
dependiendo de procesos manuales o documentos externos.

\subsection{Gestores de entrega simples}

Las soluciones m\'as ligeras se centran en el acto de subir archivos,
pero no ofrecen trazabilidad de versiones, validaci\'on avanzada
ni mecanismos de control contextual por curso y grupo. En estos
escenarios el docente debe reconstruir manualmente qu\'e se entreg\'o,
cu\'ando y bajo qu\'e reglas, lo que incrementa errores y retrabajo.

\subsection{Brechas sectoriales detectadas}

De la revisi\'on anterior se desprenden brechas recurrentes. Falta validaci\'on previa de estructura y nombre en paquetes ZIP, la trazabilidad de versiones y reenv\'ios es limitada dentro del ciclo de entrega, y las validaciones por tipo de entrega resultan insuficientes para escenarios heterog\'eneos. Adem\'as, el control contextual por curso y grupo es poco expl\'icito en herramientas generalistas, y la operativa docente sigue dependiendo de descarga y comprobaci\'on manual de lotes de entrega. Estas carencias determinan el foco de la propuesta.

\subsection{Implicaciones para el dise\~no del TFG}

Estas brechas guiaron la definici\'on del alcance funcional y de las
validaciones del sistema. La aportaci\'on se orient\'o a reducir carga
manual, mejorar la trazabilidad y ofrecer verificaciones t\'ecnicas
que no suelen estar disponibles en soluciones generalistas.

\section{Aportaci\'on realizada (escenario \emph{objetivo})}

La propuesta implementada aporta una mejora estructural,
orientada a reducir fricci\'on operativa,
aumentar trazabilidad y consolidar calidad de producto.

\subsection{Aportaci\'on funcional}

La aportaci\'on funcional responde a las brechas operativas y sectoriales
identificadas en secciones previas. En concreto, se implementaron capacidades de gesti\'on integral de cursos, grupos, actividades y entregables con visibilidad por rol, as\'i como validaci\'on avanzada de ZIP (estructura, extensiones y nombre esperado). Adem\'as, se incorpor\'o un modo estricto o m\'inimo requerido en entregas ZIP seg\'un criterio docente, control de tipos de entrega (archivo, enlace o solo texto) con reglas espec\'ificas, y versionado con reenv\'io y estado de entrega trazables. Como soporte operativo, se a\~nadi\'o descarga masiva, previsualizaci\'on de contenido para optimizar revisi\'on y un borrador autom\'atico en cliente con revalidaci\'on previa al env\'io.

Estas capacidades se alinean con los procesos objetivo del an\'alisis:
\textbf{PRON-004} (creaci\'on/publicaci\'on de actividad),
\textbf{PRON-005} (realizaci\'on de entrega con validaci\'on)
y \textbf{PRON-006} (evaluaci\'on y feedback).

\subsection{Aportaci\'on arquitect\'onica y tecnol\'ogica}

Desde la perspectiva t\'ecnica, la contribuci\'on principal se concreta en una separaci\'on frontend-backend mediante API REST, una arquitectura por capas en backend que a\'isla responsabilidades, el uso de DTOs para desacoplar contratos de API del modelo de persistencia y una evoluci\'on de esquema controlada por migraciones versionadas. Este conjunto reduce el acoplamiento y facilita mantenimiento evolutivo.

Este dise\~no reduce acoplamiento,
facilita mantenimiento
y habilita evoluci\'on incremental sin ruptura del n\'ucleo funcional.

\subsection{Aportaci\'on en calidad y operaci\'on}

Adem\'as de funcionalidad, el proyecto incorpora capacidades de calidad y operaci\'on mediante una estrategia de testing en capas (unitarias, integraci\'on, E2E y mutaci\'on), pipelines CI/CD para validaci\'on continua y despliegue reproducible en entorno PaaS con configuraci\'on por entorno. Estas decisiones garantizan que el producto sea verificable y operable en condiciones reales.

La aportaci\'on, por tanto, no se limita a ``que funcione'',
sino a que el sistema sea verificable, mantenible y operable.

\section{Matriz de mapeo debilidades-aportaciones}

La matriz siguiente sintetiza la relaci\'on entre debilidades detectadas
y la respuesta concreta implementada en el proyecto.

\begin{table}[ht]
\centering
\begin{tabular}{|P{2.5cm}|P{4.5cm}|P{4.9cm}|}
\hline
\textbf{Debilidad} & \textbf{Impacto inicial} & \textbf{Aportaci\'on aplicada} \\
\hline
DEB-001 & Sobrecarga manual en correcci\'on y organizaci\'on de ficheros & Flujo estructurado de entregas y consultas por contexto \\
\hline
DEB-002 & Dificultad de evoluci\'on sobre stack no modular & Arquitectura desacoplada con API y capas de servicio \\
\hline
DEB-003 & Fricci\'on para el estudiante durante la entrega & Formularios guiados, validaciones y seguimiento de estado \\
\hline
DEB-004 & Dependencia de almacenamiento local/heterog\'eneo & Integraciones configurables y estrategia de almacenamiento por entorno \\
\hline
\end{tabular}
\caption{Relaci\'on entre debilidades detectadas y aportaci\'on implementada}
\label{tab:mapeo-debilidades-aportacion}
\end{table}

\section{Comparativa sint\'etica}

Esta comparativa resume el salto entre escenario de partida y la
aportaci\'on del TFG en t\'erminos de operaci\'on y trazabilidad.

\begin{table}[ht]
\centering
\small
\begin{tabular}{|P{3.0cm}|P{5.3cm}|P{5.3cm}|}
\hline
\textbf{Aspecto} & \textbf{Situaci\'on inicial} & \textbf{Aportaci\'on del TFG} \\
\hline
Gesti\'on de entregas & Descarga masiva y organizaci\'on manual en local & Flujo de entrega estructurado con trazabilidad \\
\hline
Evaluaci\'on & Proceso parcialmente desacoplado de la entrega & Evaluaci\'on y feedback integrados en el mismo ciclo \\
\hline
Evoluci\'on t\'ecnica & Dependencia de stack monol\'itico/generalista & Arquitectura modular con API REST \\
\hline
Control de acceso & Autenticaci\'on institucional de base & Modelo de roles y validaci\'on contextual por operaci\'on \\
\hline
Calidad del desarrollo & Validaci\'on heterog\'enea seg\'un tarea & Estrategia de testing y automatizaci\'on CI/CD \\
\hline
\end{tabular}
\caption{Comparativa entre escenario de partida y propuesta desarrollada}
\label{tab:comparativa-antecedentes}
\end{table}

\section{Alineaci\'on con objetivos de negocio}

La aportaci\'on realizada responde de forma directa a los objetivos de negocio definidos en la fase de an\'alisis. Se materializa el \textbf{OBJN-001} al agilizar la gesti\'on docente y reducir tareas manuales de bajo valor, se satisface el \textbf{OBJN-002} al mejorar accesibilidad y usabilidad del proceso de entrega, se refuerza el \textbf{OBJN-003} al centralizar evaluaci\'on y retroalimentaci\'on con trazabilidad, y se cumple el \textbf{OBJN-004} al estructurar informaci\'on que facilita evoluci\'on e integraciones futuras.

\section{Limitaciones residuales y continuidad}

Aun cuando la mejora es significativa, se mantienen l\'ineas de continuidad planificadas: ampliaci\'on de funcionalidades avanzadas fuera del MVP, extensi\'on de integraciones externas seg\'un contexto institucional y refuerzo de observabilidad y m\'etricas operativas de producci\'on. Esta continuidad permite evolucionar sin comprometer la estabilidad alcanzada.

Estas limitaciones no invalidan la aportaci\'on obtenida,
sino que delimitan una hoja de ruta realista para evoluci\'on posterior.

\section{Conclusiones del cap\'itulo}

El an\'alisis de antecedentes confirma que el principal problema del dominio
no era la ausencia de plataformas docentes,
sino la falta de especializaci\'on en el flujo operativo de entregables.

La aportaci\'on de este TFG se sit\'ua precisamente en ese hueco:
transformar un proceso manual y costoso en un proceso digital trazable,
modular y orientado a evaluaci\'on continua,
con una base t\'ecnica preparada para evolucionar con menor riesgo de deuda estructural.
